% !TEX root = ../main.tex
\subsection*{4.2 事件分析设计：从原始材料到规范化研究对象}\label{page:15}
\noindent\textbf{事件分析将可追溯材料转化为能够触发产业研究的规范化对象。}这一任务包含四个连续步骤：识别材料中的时间、主体、对象和变化方向；匹配预定义的研究范围；绑定来源、日期和核验状态；生成可由图谱、公司、评分与证据模块共同读取的事件记录。系统已经实现已登记事件的调用与分析，自然语言事件解析将以相同记录为输出，持续拓展事件发现能力\cite{lhy-contracts,lhy-architecture}。

事件结构化首先解决“同一事实被不同模块如何理解”的问题。例如，一条涉及电池价格变化的材料，至少需要区分价格对象、变化幅度、时间、地区、可能受影响的电池或动力系统节点，以及来源的可信度。若这些信息只被写在段落中，图谱查询、公司筛选和证据复核将分别产生不同的语义边界。相反，登记事件将事实主张、研究范围和证据入口绑定到一个\texttt{event\_id}，使后续模块围绕同一对象工作。

\begin{center}\small 表15-1\quad 事件规范化的处理单元、输出字段与质量要求\par\vspace{1mm}
\begin{tabularx}{\linewidth}{p{0.19\linewidth}p{0.30\linewidth}X}
\toprule
处理单元 & 结构化输出 & 质量要求\\
\midrule
事实抽取 & 事件名称、类型、日期、地区、方向、变化对象与幅度 & 事实描述应能回到原始来源定位；时间与地区不能由后续模块默认补全。\\
研究范围匹配 & 受影响节点、行业标签、事件类别与优先级 & 节点应使用图谱可识别的标准ID或明确的待映射状态。\\
来源绑定 & 标题、URL、发布日期、摘录、来源等级与核验状态 & 记录来源本身与“来源支持何种事实”同时保留，避免链接成为孤立装饰。\\
版本登记 & \texttt{event\_id}、\texttt{data\_version}与创建/修订信息 & 已登记版本可被复跑；后续修订形成新记录或清晰的版本关联。\\
任务触发 & 图谱查询条件、证据检索键、公司筛选范围和运行请求 & 下游模块读取事件字段，而不重新解释输入文本。\\
\bottomrule
\end{tabularx}\end{center}

在当前版本中，事件由研究者从已有列表选择并提交，系统据此发起分析请求。这一范围使演示链路可重复，也为事件对象、路径对象和报告对象之间的接口提供稳定测试条件。随着事件库和证据库接入，系统将基于统一字段扩展自然语言抽取、时空范围识别和候选关系发现；这些模型输出将被标记为候选或待核验记录，并经来源、字段和人工复核后进入正式分析。

\noindent\textbf{事件规范化让“信息出现”立即转化为可启动、可复核的研究任务。}政策文本、公告、价格材料和行业报告经由统一事件对象进入主链，产业路径、公司比较和财务解释共享同一个事实起点，从而使全链报告具备一致的可追溯性。
